\chapter{组级最优通行带与最大间隙策略}

本章对应MIKU的第二项改进，针对Apollo PathBoundsDecider对每个障碍物独立决定避让方向所导致的矛盾决策与路径BLOCKED问题，提出基于扫描线分组与最大间隙策略的组级全局最优算法。本章先分析Apollo \texttt{UpdatePathBoundary\-BySLPolygon}函数的逐障碍物贪心决策逻辑及其在两侧夹击场景下的失败模式，再构建基于纵向重叠的连通分量分组方法，在组内证明最优分割连续性引理与最大间隙最优性定理，将组合搜索复杂度从$O(2^k)$降至$O(k\log k)$。本章使用的差异化安全裕度$\delta_i$来自第四章的多因子威胁度评估，最终输出的组级避让方向决策与通行带宽度将作为第六章ST走廊边界构造的SL层基础。

\section{Apollo 逐障碍物决策的源码缺陷与失败模式}
\label{sec:problem_sl}

\subsection{源码逻辑与决策入口}

PathBoundsDecider是Apollo路径规划阶段的关键模块，负责在每个离散纵向位置$s$处构建横向路径边界$[l^-(s),\; l^+(s)]$。问题在于它对每个障碍物独立地决定避让方向（LEFT\_NUDGE或RIGHT\_NUDGE），而不考虑同一$s$截面上其他活跃障碍物的存在。当道路两侧同时存在障碍物时，这种逐障碍物的贪心决策可能产生相互矛盾的避让方向，即左侧障碍物使上边界下压、右侧障碍物使下边界上抬，导致$l^+(s) < l^-(s)$，路径边界被判定为BLOCKED，最终触发FallbackPath停车。

问题根源在\texttt{UpdatePathBoundary\-BySLPolygon}函数，其逻辑分三步：逐障碍物独立决定避让方向、根据避让方向收缩路径边界、检测边界交叉并截断路径。代码\ref{lst:update_path_boundary}展示了该函数的关键流程。

\begin{lstlisting}[language=C++,caption={UpdatePathBoundaryBySLPolygon核心逻辑},label={lst:update_path_boundary}]
// path_bounds_decider_util.cc: UpdatePathBoundaryBySLPolygon
for (size_t i = 1; i < path_boundary->size(); ++i) {
  auto& left_bound = path_boundary->at(i).l_upper;   // 上边界
  auto& right_bound = path_boundary->at(i).l_lower;  // 下边界
  for (size_t j = 0; j < sl_polygon->size(); j++) {
    // --- 步骤1: 逐障碍物独立决定避让方向 ---
    if (sl_polygon->at(j).NudgeInfo() == SLPolygon::UNDEFINED) {
      double obs_l = (l_lower + l_upper) / 2;  // 障碍物中心横向位置
      if (is_passable_right && is_passable_left) {
        // 两侧均可通行时，根据障碍物中心与参考线的相对位置决定
        if (last_max_nudge_l < obs_l) {
          sl_polygon->at(j).SetNudgeInfo(SLPolygon::RIGHT_NUDGE);
        } else {
          sl_polygon->at(j).SetNudgeInfo(SLPolygon::LEFT_NUDGE);
        }
      } else if (is_passable_right) {
        sl_polygon->at(j).SetNudgeInfo(SLPolygon::RIGHT_NUDGE);
      } else {
        sl_polygon->at(j).SetNudgeInfo(SLPolygon::LEFT_NUDGE);
      }
    }
    // --- 步骤2: 根据避让方向更新边界 ---
    if (sl_polygon->at(j).NudgeInfo() == SLPolygon::RIGHT_NUDGE) {
      if (!sl_polygon->at(j).is_passable()[RIGHT_INDEX]) {
        sl_polygon->at(j).SetNudgeInfo(SLPolygon::BLOCKED);
        break;  // 右侧不可通行，标记BLOCKED
      }
      // 上边界收缩至障碍物左边缘减安全距离
      if (obs_right_nudge_bound.l_upper.l < left_bound.l) {
        left_bound.l = obs_right_nudge_bound.l_upper.l;
        left_bound.type = BoundType::OBSTACLE;
      }
    } else if (sl_polygon->at(j).NudgeInfo() == SLPolygon::LEFT_NUDGE) {
      if (!sl_polygon->at(j).is_passable()[LEFT_INDEX]) {
        sl_polygon->at(j).SetNudgeInfo(SLPolygon::BLOCKED);
        break;  // 左侧不可通行，标记BLOCKED
      }
      // 下边界上抬至障碍物右边缘加安全距离
      if (obs_left_nudge_bound.l_lower.l > right_bound.l) {
        right_bound.l = obs_left_nudge_bound.l_lower.l;
        right_bound.type = BoundType::OBSTACLE;
      }
    }
  }
  // --- 步骤3: BLOCKED判定与路径截断 ---
  if (!blocked_id->empty()) {
    TrimPathBounds(i, path_boundary);  // 截断后续路径边界
    return false;                      // 触发FallbackPath停车
  }
}
\end{lstlisting}

步骤1对每个在位置$s$处活跃的障碍物，通过\texttt{SetNudgeInfo}函数根据障碍物中心相对于参考线的横向位置独立决定\texttt{RIGHT\_NUDGE}即从右侧绕行使上边界下压，或\texttt{LEFT\_NUDGE}即从左侧绕行使下边界上抬。步骤2根据避让方向对路径边界进行收缩，安全距离\texttt{delta}由自车半宽即\texttt{GetBufferBetweenADCCenterAndEdge}的返回值与额外缓冲参数\texttt{obstacle\_lat\_buffer}相加得到。步骤3在上下边界交叉即存在BLOCKED标记的障碍物时，调用\texttt{TrimPathBounds}清除后续路径边界并返回\texttt{false}，最终触发\texttt{FallbackPath}生成减速停车轨迹。

\subsection{独立决策的失败模式}

以下通过一个对称场景说明独立决策如何导致不必要的停车。

\textbf{场景设定}\quad 双车道道路总宽7.5\,m，$l_{road}^- = -3.75$\,m，$l_{road}^+ = 3.75$\,m。在某纵向位置$s^*$处同时存在两个路锥，关于参考线近似对称

\begin{itemize}[noitemsep]
\item $O_a$偏右，横向占据$[0.1,\; 0.4]$\,m，中心$0.25$\,m
\item $O_b$偏左，横向占据$[-0.4,\; -0.1]$\,m，中心$-0.25$\,m
\end{itemize}

两个路锥均未越过车道中线，安全距离$\delta = 1.2$\,m。

\textbf{Apollo独立决策}\quad $O_a$中心$0.25 > 0$，分配RIGHT\_NUDGE，上边界收缩至
\[
l^+(s^*) = 0.1 - 1.2 = -1.1\,\text{m}
\]
$O_b$中心$-0.25 < 0$，分配LEFT\_NUDGE，下边界抬升至
\[
l^-(s^*) = -0.1 + 1.2 = 1.1\,\text{m}
\]
$l^+ = -1.1 < l^- = 1.1$，判定BLOCKED。贪心策略将自车挤入两个路锥之间仅0.2\,m的间隙，路径不可行，触发停车。

\textbf{全局统一决策}\quad 将两个路锥作为一组，统一分配RIGHT\_NUDGE，自车从两者右侧绕行
\[
l^+(s^*) = \min(0.1 - 1.2,\; -0.4 - 1.2) = -1.6\,\text{m}, \quad l^-(s^*) = -3.75\,\text{m}
\]
通行带宽度$W = -1.6 - (-3.75) = 2.15$\,m，自车可借用相邻车道从右侧安全通过。由对称性，统一LEFT\_NUDGE同样可行，$W = 3.75 - 1.6 = 2.15$\,m。

贪心策略将两个偏向不同侧的障碍物分别向内挤压，全局协调只需统一绕行方向，即可利用道路另一侧的充裕空间。当$k$个障碍物在同一纵向位置活跃时，避让方向共有$2^k$种组合，贪心策略只检查了其中一种，可能恰好选中不可行的组合而遗漏可行解。该问题的形式化建模在第三章统一框架下已经给出，高效求解方法将在本章下一节详细展开。

图\ref{fig:path_conservative}展示了该场景中独立决策与全局协调的对比。

\begin{figure}[htbp]
\centering
\begin{subfigure}[b]{0.48\textwidth}
  \centering
  \resizebox{\linewidth}{!}{% 子图(a)：独立决策，矛盾避让 → BLOCKED
\begin{tikzpicture}[scale=0.85, transform shape]
  % 道路背景
  \fill[bglight] (-0.5,-2.5) rectangle (8.5,2.5);
  \draw[gray, thick] (-0.5, 2.5) -- (8.5, 2.5);
  \draw[gray, thick] (-0.5,-2.5) -- (8.5,-2.5);
  \draw[dashed, gray!60] (-0.5, 0) -- (8.5, 0);
  \node[gray, font=\scriptsize, anchor=east] at (-0.6, 0) {$l{=}0$};

  % O_a（上方）
  \fill[dangerred!70, draw=dangerred, thick] (3.2, 0.4) rectangle (4.8, 0.8);
  \node[font=\scriptsize, white] at (4.0, 0.6) {$O_a$};

  % O_b（下方）
  \fill[dangerred!70, draw=dangerred, thick] (3.2,-0.8) rectangle (4.8,-0.4);
  \node[font=\scriptsize, white] at (4.0,-0.6) {$O_b$};

  % 膨胀虚线 δ=0.6
  \draw[dangerred!60, dashed, thick] (2.6,-0.2) rectangle (5.4, 1.4);
  \node[dangerred!60, font=\scriptsize, anchor=west] at (5.5, 1.1) {$O_a{+}\delta$};
  \draw[dangerred!60, dashed, thick] (2.6,-1.4) rectangle (5.4, 0.2);
  \node[dangerred!60, font=\scriptsize, anchor=west] at (5.5,-1.1) {$O_b{+}\delta$};

  % 上界 l+(s)：右绕 O_a
  \draw[deeppink, very thick]
    (-0.5, 2.5) -- (2.4, 2.5) -- (2.4, -0.3) -- (5.6, -0.3) -- (5.6, 2.5) -- (8.5, 2.5);
  \node[deeppink, font=\scriptsize] at (7.5, 2.1) {$l^+(s)$};

  % 下界 l-(s)：左绕 O_b
  \draw[deepblue, very thick]
    (-0.5,-2.5) -- (2.4,-2.5) -- (2.4, 0.3) -- (5.6, 0.3) -- (5.6,-2.5) -- (8.5,-2.5);
  \node[deepblue, font=\scriptsize] at (7.5,-2.1) {$l^-(s)$};

  % 矛盾标注
  \draw[dangerred, thick, {Stealth}-{Stealth}] (6.5,-0.3) -- (6.5, 0.3);
  \node[dangerred, font=\scriptsize, anchor=west] at (6.6, 0) {$l^+ < l^-$};

  % BLOCKED 判定
  \node[draw=dangerred, fill=dangerred!15, font=\scriptsize,
        rounded corners=3pt, align=center] at (4.0,-2.1)
    {$l^+ < l^-$ $\Rightarrow$ \textbf{BLOCKED}};
\end{tikzpicture}
}
  \caption{独立决策，矛盾避让 $\to$ BLOCKED}
  \label{fig:path_conservative:a}
\end{subfigure}\hfill
\begin{subfigure}[b]{0.48\textwidth}
  \centering
  \resizebox{\linewidth}{!}{% 子图(b)：全局协调，统一右绕 → 通行
\begin{tikzpicture}[scale=0.85, transform shape]
  % 道路背景
  \fill[bglight] (-0.5,-2.5) rectangle (8.5,2.5);
  \draw[gray, thick] (-0.5, 2.5) -- (8.5, 2.5);
  \draw[gray, thick] (-0.5,-2.5) -- (8.5,-2.5);
  \draw[dashed, gray!60] (-0.5, 0) -- (8.5, 0);
  \node[gray, font=\scriptsize, anchor=east] at (-0.6, 0) {$l{=}0$};

  % O_a（上方）
  \fill[dangerred!70, draw=dangerred, thick] (3.2, 0.4) rectangle (4.8, 0.8);
  \node[font=\scriptsize, white] at (4.0, 0.6) {$O_a$};

  % O_b（下方）
  \fill[dangerred!70, draw=dangerred, thick] (3.2,-0.8) rectangle (4.8,-0.4);
  \node[font=\scriptsize, white] at (4.0,-0.6) {$O_b$};

  % 膨胀虚线 δ=0.6
  \draw[dangerred!60, dashed, thick] (2.6,-0.2) rectangle (5.4, 1.4);
  \node[dangerred!60, font=\scriptsize, anchor=west] at (5.5, 0.6) {$O_a{+}\delta$};
  \draw[dangerred!60, dashed, thick] (2.6,-1.4) rectangle (5.4, 0.2);
  \node[dangerred!60, font=\scriptsize, anchor=west] at (5.5,-0.6) {$O_b{+}\delta$};

  % 上界 l+(s)：统一右绕，压到 O_b 膨胀下缘
  \draw[deeppink, very thick]
    (-0.5, 2.5) -- (2.4, 2.5) -- (2.4, -1.5) -- (5.6, -1.5) -- (5.6, 2.5) -- (8.5, 2.5);
  \node[deeppink, font=\scriptsize] at (7.5, 2.1) {$l^+(s)$};

  % 下界 l-(s)：道路下界
  \draw[deepblue, very thick, dashed]
    (-0.5,-2.5) -- (8.5,-2.5);
  \node[deepblue, font=\scriptsize, left] at (-0.6,-2.5) {$l^-(s)$};

  % 通行带
  \fill[lightgreen, opacity=0.25] (-0.5,-2.5) rectangle (8.5,-1.5);

  % 通行带宽度标注
  \draw[deeppink, thick, {Stealth}-{Stealth}] (0.5,-2.5) -- (0.5,-1.5);
  \node[deeppink, font=\scriptsize, anchor=east] at (0.4,-2.0) {$W{=}1.0$m};

  % 自车轨迹
  \draw[deepblue, very thick, -{Stealth}]
    plot[smooth, tension=0.55] coordinates
      {(0.5, 0) (1.5,-0.5) (2.5,-1.8) (4.0,-2.0) (5.5,-1.8) (6.5,-0.5) (7.5, 0)};
  \node[deepblue, font=\scriptsize] at (4.0,-2.45) {规划轨迹};

  % PASS 判定
  \node[draw=lightgreen!80!black, fill=lightgreen!20, font=\scriptsize,
        rounded corners=3pt, align=center] at (4.0, 2.0)
    {统一右绕 $\Rightarrow$ \textbf{PASS}};
\end{tikzpicture}
}
  \caption{全局协调，统一右绕 $\to$ 通行}
  \label{fig:path_conservative:b}
\end{subfigure}
\caption{独立决策导致矛盾避让方向与全局协调的对比}
\label{fig:path_conservative}
\end{figure}

除了两个障碍物的极端情形外，随着障碍物数量增加，独立决策导致的可行域压缩会以近似指数方式恶化。图\ref{fig:feasible_compress}给出了同一车道内从2个、4个到极端密度障碍物场景下的可行域演化。2个障碍物场景下，尚存较宽的可行通带可供自车选择，尽管贪心决策已开始压缩通带宽度；4个障碍物场景下，多次独立收缩使可行域退化为断续的狭窄通带，任一处出现贪心矛盾即触发BLOCKED；极端密度场景下，可行域进一步碎片化，几乎任何贪心决策组合都会导致某处边界失效。可见在高密度场景下，改进决策粒度对可行性恢复的贡献度远高于单纯调整安全裕度。

\begin{figure}[htbp]
\centering
\begin{subfigure}[b]{0.32\textwidth}
  \centering
  \resizebox{\linewidth}{!}{% 2 个障碍物 通行带可行（子图 a）
\begin{tikzpicture}[scale=0.7, transform shape]
  \draw[-Stealth, thick] (0,-2.5) -- (8.5,-2.5) node[right, font=\footnotesize] {$s$};
  \draw[-Stealth, thick] (0,-2.5) -- (0,2.5) node[above, font=\footnotesize] {$l$};

  \draw[gray, thick] (0,1.875) -- (8,1.875);
  \draw[gray, thick] (0,-1.875) -- (8,-1.875);
  \draw[gray, dashed] (0,0) -- (8,0);

  \begin{scope}[on background layer]
    \fill[bglight] (0,-1.875) rectangle (8,1.875);
  \end{scope}

  % O1（上方靠边）y=[0.6, 1.3], s=[2, 4]
  \fill[dangerred!70, draw=dangerred, thick] (2,0.6) rectangle (4,1.3);
  \node[white, font=\footnotesize\bfseries] at (3,0.95) {$O_1$};

  % O2（下方靠边）y=[-1.2, -0.5], s=[5, 7]
  \fill[dangerred!70, draw=dangerred, thick] (5,-1.2) rectangle (7,-0.5);
  \node[white, font=\footnotesize\bfseries] at (6,-0.85) {$O_2$};

  % O1 右绕 → 上界压到 O1 膨胀下缘 (0.6-0.4=0.2)
  \draw[deeppink, very thick]
    (0,1.875) -- (1.8,1.875) -- (1.8,0.2) -- (4.2,0.2) -- (4.2,1.875) -- (8,1.875);
  \node[deeppink, font=\tiny] at (1, 2.1) {$l^+(s)$};

  % O2 左绕 → 下界抬到 O2 膨胀上缘 (-0.5+0.4=-0.1)
  \draw[deepblue, very thick]
    (0,-1.875) -- (4.8,-1.875) -- (4.8,-0.1) -- (7.2,-0.1) -- (7.2,-1.875) -- (8,-1.875);
  \node[deepblue, font=\tiny] at (7.5, -2.1) {$l^-(s)$};

  \node[lightgreen!60!black, font=\footnotesize\bfseries] at (4,-1.35) {通行带充足};
\end{tikzpicture}
}
  \caption{2 个障碍物}
  \label{fig:feasible_compress:a}
\end{subfigure}\hfill
\begin{subfigure}[b]{0.32\textwidth}
  \centering
  \resizebox{\linewidth}{!}{% 4 个障碍物 通行带狭窄（子图 b）
\begin{tikzpicture}[scale=0.7, transform shape]
  \draw[-Stealth, thick] (0,-2.5) -- (8.5,-2.5) node[right, font=\footnotesize] {$s$};
  \draw[-Stealth, thick] (0,-2.5) -- (0,2.5) node[above, font=\footnotesize] {$l$};

  \draw[gray, thick] (0,1.875) -- (8,1.875);
  \draw[gray, thick] (0,-1.875) -- (8,-1.875);
  \draw[gray, dashed] (0,0) -- (8,0);

  \begin{scope}[on background layer]
    \fill[bglight] (0,-1.875) rectangle (8,1.875);
  \end{scope}

  % O1（上方靠边）y=[0.5, 1.1], s=[1, 2.5]
  \fill[dangerred!70, draw=dangerred, thick] (1,0.5) rectangle (2.5,1.1);
  \node[white, font=\tiny\bfseries] at (1.75,0.8) {$O_1$};

  % O2（下方靠边）y=[-1.0, -0.4], s=[2.5, 4]
  \fill[dangerred!70, draw=dangerred, thick] (2.5,-1.0) rectangle (4,-0.4);
  \node[white, font=\tiny\bfseries] at (3.25,-0.7) {$O_2$};

  % O3（上方靠边）y=[0.4, 0.9], s=[4.5, 6]
  \fill[dangerred!70, draw=dangerred, thick] (4.5,0.4) rectangle (6,0.9);
  \node[white, font=\tiny\bfseries] at (5.25,0.65) {$O_3$};

  % O4（下方靠边）y=[-0.8, -0.2], s=[6, 7.5]
  \fill[dangerred!70, draw=dangerred, thick] (6,-0.8) rectangle (7.5,-0.2);
  \node[white, font=\tiny\bfseries] at (6.75,-0.5) {$O_4$};

  % 上界 l+(s)：O1右绕压到0.1, O3右绕压到0.0
  \draw[deeppink, very thick]
    (0,1.875) -- (0.8,1.875) -- (0.8,0.1) -- (2.7,0.1)
    -- (2.7,1.875) -- (4.3,1.875) -- (4.3,0.0) -- (6.2,0.0)
    -- (6.2,1.875) -- (8,1.875);

  % 下界 l-(s)：O2左绕抬到0.0, O4左绕抬到0.2
  \draw[deepblue, very thick]
    (0,-1.875) -- (2.3,-1.875) -- (2.3,0.0) -- (4.2,0.0)
    -- (4.2,-1.875) -- (5.8,-1.875) -- (5.8,0.2) -- (7.7,0.2)
    -- (7.7,-1.875) -- (8,-1.875);

  % 窄处标注
  \draw[lightorange, ultra thick] (2.5,0.0) -- (2.5,0.1);
  \node[lightorange!80!black, font=\tiny\bfseries, anchor=west] at (2.6,0.05) {窄};
  \draw[lightorange, ultra thick] (6,0.0) -- (6,0.2);
  \node[lightorange!80!black, font=\tiny\bfseries, anchor=west] at (6.1,0.1) {窄};

  \node[deeppink, font=\tiny] at (0.5, 2.1) {$l^+(s)$};
  \node[deepblue, font=\tiny] at (7.5, -2.1) {$l^-(s)$};
  \node[lightorange!80!black, font=\footnotesize\bfseries] at (4,-1.5) {多次收缩 通行带退化};
\end{tikzpicture}
}
  \caption{4 个障碍物}
  \label{fig:feasible_compress:b}
\end{subfigure}\hfill
\begin{subfigure}[b]{0.32\textwidth}
  \centering
  \resizebox{\linewidth}{!}{% 密集障碍物 BLOCKED（子图 c）
\begin{tikzpicture}[scale=0.7, transform shape]
  \draw[-Stealth, thick] (0,-2.5) -- (8.5,-2.5) node[right, font=\footnotesize] {$s$};
  \draw[-Stealth, thick] (0,-2.5) -- (0,2.5) node[above, font=\footnotesize] {$l$};

  \draw[gray, thick] (0,1.875) -- (8,1.875);
  \draw[gray, thick] (0,-1.875) -- (8,-1.875);
  \draw[gray, dashed] (0,0) -- (8,0);

  \begin{scope}[on background layer]
    \fill[bglight] (0,-1.875) rectangle (8,1.875);
  \end{scope}

  % 密集障碍物（交替上下，靠近道路边缘）
  \fill[dangerred!70, draw=dangerred, thick] (0.5,0.4) rectangle (1.5,1.2);
  \node[white, font=\tiny\bfseries] at (1,0.8) {$O_1$};

  \fill[dangerred!70, draw=dangerred, thick] (1.5,-1.1) rectangle (2.5,-0.3);
  \node[white, font=\tiny\bfseries] at (2,-0.7) {$O_2$};

  \fill[dangerred!70, draw=dangerred, thick] (2.5,0.3) rectangle (3.5,1.1);
  \node[white, font=\tiny\bfseries] at (3,0.7) {$O_3$};

  \fill[dangerred!70, draw=dangerred, thick] (3.5,-1.0) rectangle (4.5,-0.2);
  \node[white, font=\tiny\bfseries] at (4,-0.6) {$O_4$};

  \fill[dangerred!70, draw=dangerred, thick] (4.5,0.5) rectangle (5.5,1.3);
  \node[white, font=\tiny\bfseries] at (5,0.9) {$O_5$};

  \fill[dangerred!70, draw=dangerred, thick] (5.5,-1.2) rectangle (6.5,-0.3);
  \node[white, font=\tiny\bfseries] at (6,-0.75) {$O_6$};

  \fill[dangerred!70, draw=dangerred, thick] (6.5,0.3) rectangle (7.5,1.1);
  \node[white, font=\tiny\bfseries] at (7,0.7) {$O_7$};

  % 上界 l+(s)：逐步压低 (δ=0.4)
  \draw[deeppink, very thick]
    (0,1.875) -- (0.3,1.875) -- (0.3,0.0) -- (1.7,0.0)
    -- (1.7,1.875) -- (2.3,1.875) -- (2.3,-0.1) -- (3.7,-0.1)
    -- (3.7,1.875) -- (4.3,1.875) -- (4.3,0.1) -- (5.7,0.1);

  % 下界 l-(s)：逐步抬高
  \draw[deepblue, very thick]
    (0,-1.875) -- (1.3,-1.875) -- (1.3,0.1) -- (2.7,0.1)
    -- (2.7,-1.875) -- (3.3,-1.875) -- (3.3,0.2) -- (4.7,0.2);

  % BLOCKED 点
  \fill[dangerred] (4.5,0.15) circle (3pt);

  % 截断区域
  \fill[dangerred, opacity=0.1] (4.5,-1.875) rectangle (8,1.875);
  \node[dangerred, font=\footnotesize\bfseries] at (6.5,1.3) {$l^+ < l^- \rightarrow$ BLOCKED};
  \node[dangerred, font=\tiny] at (6.5,-1.5) {TrimPathBounds 截断};

  \node[deeppink, font=\tiny] at (0.5, 2.1) {$l^+(s)$};
  \node[deepblue, font=\tiny] at (0.5, -2.1) {$l^-(s)$};
\end{tikzpicture}
}
  \caption{密集障碍物}
  \label{fig:feasible_compress:c}
\end{subfigure}
\caption{多障碍物场景下可行域的退化演化}
\label{fig:feasible_compress}
\end{figure}

\subsection{BLOCKED 后的截断与恢复机制局限}

当PathBoundsDecider判定BLOCKED后，调用\texttt{TrimPathBounds}截断路径

\begin{lstlisting}[language=C++,caption={TrimPathBounds路径截断},label={lst:trim_path}]
// path_bounds_decider_util.cc: TrimPathBounds
void PathBoundsDeciderUtil::TrimPathBounds(
    const int path_blocked_idx, PathBoundary* const path_boundaries) {
  if (path_blocked_idx != -1) {
    if (path_blocked_idx == 0) {
      AINFO << "Completely blocked. Cannot move at all.";
    }
    double front_edge_to_center =
        VehicleConfigHelper::GetConfig().vehicle_param().front_edge_to_center();
    double trimmed_s =
        path_boundaries->at(path_blocked_idx).s - front_edge_to_center;
    while (path_boundaries->size() > 1 &&
           path_boundaries->back().s > trimmed_s) {
      path_boundaries->pop_back();  // 逐个弹出BLOCKED位置之后的边界点
    }
  }
}
\end{lstlisting}

该函数将路径边界截断至BLOCKED位置的前方即减去车头到中心的距离，后续\texttt{FallbackPath}模块基于截断后的短路径生成减速停车轨迹。一旦某个避让方向组合导致BLOCKED，系统不会回溯尝试其他方向组合，而是直接进入停车流程。

Apollo提供的\texttt{LaneBorrowPath}机制可以在一定条件下缓解BLOCKED问题，但该机制的触发条件非常严格。代码\ref{lst:lane_borrow_conditions}展示了借道决策需要通过的5个前置检查及条件4的具体实现

\begin{lstlisting}[language=C++,caption={LaneBorrowPath的前置条件与长期阻塞判定},label={lst:lane_borrow_conditions}]
// lane_borrow_path.cc: IsNecessaryToBorrowLane
// 条件1: 仅单参考线场景（路口等多参考线场景不借道）
if (!HasSingleReferenceLine(*frame_)) { return false; }
// 条件2: 自车速度低于阈值（默认5.0 m/s）
if (!IsWithinSidePassingSpeedADC(*frame_)) { return false; }
// 条件3: 阻塞障碍物远离路口
if (!IsBlockingObstacleFarFromIntersection(*ref_info)) { return false; }
// 条件4: 障碍物持续阻塞达到cycle阈值（默认3个周期）
if (!IsLongTermBlockingObstacle()) { return false; }
// 条件5: 阻塞障碍物在目的地范围内
if (!IsBlockingObstacleWithinDestination(*ref_info)) { return false; }

// 条件4的具体实现: IsLongTermBlockingObstacle
bool LaneBorrowPath::IsLongTermBlockingObstacle() {
  if (injector_->planning_context()->planning_status().path_decider()
          .front_static_obstacle_cycle_counter() >=
      config_.long_term_blocking_obstacle_cycle_threshold()) {  // 默认值: 3
    return true;
  }
  return false;
}
\end{lstlisting}

即使存在可行的借道方案，系统也必须先经历至少3个规划周期约300\,ms的停车等待才会尝试借道，且借道仅在单参考线、低速、远离路口等条件同时满足时才被允许。在路口附近、高速行驶或多参考线场景下，LaneBorrow机制不会触发，BLOCKED判定将直接导致停车。
\section{纵向重叠分组与连通分量}
\label{sec:grouping_connected}
本节起所有$u_i$和$v_i$均使用式\eqref{eq:def_u_dynamic}--\eqref{eq:def_v_dynamic}中的动态安全距离$\delta_i$计算，后续引理和定理的证明不依赖$\delta$是否为常数，仅依赖$u$和$v$值的排序性。
在全路径范围内，不同纵向位置$s$处的活跃障碍物集合不同，因此式\eqref{eq:joint_opt}的优化不能在单一截面上完成。关键观察是，如果两组障碍物在纵向上完全不重叠，则它们的避让决策相互独立。

定义纵向重叠关系。两个障碍物$O_i$和$O_j$纵向重叠，当且仅当它们的纵向区间有交集

\begin{equation}
    O_i \sim O_j \iff [s_i^{-}, s_i^{+}] \cap [s_j^{-}, s_j^{+}] \neq \emptyset
    \label{eq:overlap_relation}
\end{equation}

该重叠关系是对称的但不具有传递性。取其传递闭包，即存在一系列中间障碍物依次纵向重叠，将$n$个障碍物划分为$m$个连通分量$C_1, C_2, \ldots, C_m$。同一连通分量内的障碍物通过重叠链相连，不同连通分量之间在纵向上完全分离。由于不同连通分量在纵向上不重叠，其避让决策互不影响，可以独立求解。

连通分量的检测采用扫描线算法，将所有障碍物按$s_i^{-}$升序排列，维护一个变量$s_{max}$记录当前连通分量中所有障碍物的最大$s^{+}$值。扫描过程中，若下一个障碍物的$s^{-}$不超过$s_{max}$，则将其归入当前分量并更新$s_{max}$；否则开始新的分量。该过程的时间复杂度为$O(n \log n)$其中排序占主导，扫描本身仅需$O(n)$。

纵向上彼此重叠的障碍物需要在同一个横向截面上协调决策，因此必须归为同一组；纵向上完全分离的障碍物组，自车在通过一组时尚未遇到下一组，决策互不影响。问题因此天然分解为多个独立子问题，这是算法高效性的基础。图\ref{fig:scanline_grouping}以六障碍物算例展示了扫描过程，$s_{max}$扫描线沿$s$轴从左到右推进，遇到与当前组重叠的障碍物即并入，否则切出新分量，最终得到三个互不重叠的连通分量$C_1$到$C_3$。

\begin{figure}[htbp]
\centering
\resizebox{\textwidth}{!}{% 扫描线按 s_i^- 升序合并连通分量的过程
\begin{tikzpicture}[scale=0.95, transform shape,
    obs/.style={fill=dangerred!70, draw=dangerred, thick, rounded corners=1pt},
    grp/.style={draw=deeppink, very thick, dashed, rounded corners=3pt},
]

% s 轴
\draw[-Stealth, thick] (0,0) -- (14.5,0) node[right, font=\small] {$s$};
\foreach \x/\lab in {0/0, 2/10, 4/20, 6/30, 8/40, 10/50, 12/60, 14/70}
    \draw[thick] (\x,-0.1) -- (\x,0.1) node[below=0.1cm, font=\scriptsize] {\lab};

% 六个障碍物的 [s^-, s^+] 区间（按 s^- 升序）
% O1: [5, 15]  O2: [10, 22]  O3: [30, 38]  O4: [34, 44]  O5: [42, 50]  O6: [60, 68]
\fill[obs] (1,0.6) rectangle (3,1.1);
\node[white, font=\scriptsize\bfseries] at (2,0.85) {$O_1$};

\fill[obs] (2,1.4) rectangle (4.4,1.9);
\node[white, font=\scriptsize\bfseries] at (3.2,1.65) {$O_2$};

\fill[obs] (6,0.6) rectangle (7.6,1.1);
\node[white, font=\scriptsize\bfseries] at (6.8,0.85) {$O_3$};

\fill[obs] (6.8,1.4) rectangle (8.8,1.9);
\node[white, font=\scriptsize\bfseries] at (7.8,1.65) {$O_4$};

\fill[obs] (8.4,0.6) rectangle (10,1.1);
\node[white, font=\scriptsize\bfseries] at (9.2,0.85) {$O_5$};

\fill[obs] (12,0.6) rectangle (13.6,1.1);
\node[white, font=\scriptsize\bfseries] at (12.8,0.85) {$O_6$};

% s_max 扫描线（当前最大上界）
\draw[deepblue, very thick, densely dotted] (4.4,-0.3) -- (4.4,3.0);
\node[deepblue, font=\scriptsize, anchor=south] at (4.4,3.0) {$s_{\max}=22$};

\draw[deepblue, very thick, densely dotted] (10,-0.3) -- (10,3.0);
\node[deepblue, font=\scriptsize, anchor=south] at (10,3.0) {$s_{\max}=50$};

\draw[deepblue, very thick, densely dotted] (13.6,-0.3) -- (13.6,3.0);
\node[deepblue, font=\scriptsize, anchor=south] at (13.6,3.0) {$s_{\max}=68$};

% 三个连通分量分组框
\draw[grp] (0.7,0.4) rectangle (4.7,2.1);
\node[deeppink, font=\small\bfseries, anchor=south] at (2.7,2.1) {$C_1=\{O_1,O_2\}$};
\node[deeppink!80, font=\scriptsize, anchor=north west] at (0.78,0.43) {$s_{\min}{=}5$};

\draw[grp] (5.7,0.4) rectangle (10.3,2.1);
\node[deeppink, font=\small\bfseries, anchor=south] at (8,2.1) {$C_2=\{O_3,O_4,O_5\}$};
\node[deeppink!80, font=\scriptsize, anchor=north west] at (5.78,0.43) {$s_{\min}{=}30$};

\draw[grp] (11.7,0.4) rectangle (13.9,2.1);
\node[deeppink, font=\small\bfseries, anchor=south] at (12.8,2.1) {$C_3=\{O_6\}$};
\node[deeppink!80, font=\scriptsize, anchor=north west] at (11.78,0.43) {$s_{\min}{=}60$};

\end{tikzpicture}
}
\caption{扫描线算法按$s_i^-$升序合并连通分量的过程}
\label{fig:scanline_grouping}
\end{figure}

当障碍物仅在纵向上部分重叠即存在某些$s$截面上只有部分障碍物活跃时，组内处理采用保守近似，以所有障碍物均同时活跃为假设计算间隙。这可能导致轻微的次优性，但算法的正确性可由以下两步论证保证。首先，在最密集截面上按定理1选定的分割点$p^{*}$将障碍物划分为左绕集合$\mathcal{L}_{p^{*}}$和右绕集合$\mathcal{R}_{p^{*}}$，此时通行带上下界由$l^{+} = \min_{i \in \mathcal{R}_{p^{*}}} u_i$和$l^{-} = \max_{i \in \mathcal{L}_{p^{*}}} v_i$给出。其次，任意稀疏截面对应的活跃障碍物集合均为最密集截面的子集，移除部分障碍物只会使$\min$运算作用于更大的集合从而放宽$l^{+}$的上界，同时使$\max$运算作用于更大的集合从而放宽$l^{-}$的下界。若$\mathcal{L}_{p^{*}}$或$\mathcal{R}_{p^{*}}$被移空，则相应方向退化为道路边界。因此通行带宽度$W = l^{+} - l^{-}$单调不减，最密集截面上的决策在所有稀疏截面上均保持可行，且通行带不会更窄。

全路径的最小通行带宽度由最窄的连通分量决定，即木桶效应

\begin{equation}
    W^{*} = \min_{j=1,\ldots,m} W^{*}(C_j)
    \label{eq:global_width}
\end{equation}

下面聚焦于单个连通分量的最优求解。设该分量内有$k$个障碍物，在其公共$s$截面上按$u_i$值升序排列为$O_{(1)}, O_{(2)}, \ldots, O_{(k)}$，使得$u_{(1)} \leq u_{(2)} \leq \cdots \leq u_{(k)}$。按$u_i$排序而非$l_i^{-}$排序的原因在于引入动态安全距离$\delta_i$后，$u_i = l_i^{-} - \delta_i$的排序可能不再与$l_i^{-}$一致，详见\S\ref{sec:dynamic_delta_synthesis}，而后续引理和定理的证明仅依赖$u$值的排序性。问题转化为在$k$个按$u_i$排列的障碍物上，如何分配左绕/右绕决策以最大化通行带宽度。

从几何直觉看，$k$个按横向位置排列的障碍物将车道的横向空间分割为$k+1$个通道，包括最右侧障碍物与道路右边界之间、相邻障碍物之间、以及最左侧障碍物与道路左边界之间的间隙。最优策略是选择其中最宽的通道通行，通道一侧的所有障碍物划归左绕、另一侧的划归右绕，形成一个在排序中连续的分割。搜索空间从$2^k$个任意分配降为$k+1$个候选通道的选取。以下引理和定理给出严格证明，并建立通行带宽度与间隙之间的精确等式。

\section{最优分割连续性与最大间隙定理}
本节给出MIKU通行带求解的两个核心理论结果。引理1将组合搜索空间从$2^k$降为$k+1$个连续分割，定理1把通行带宽度与候选间隙建立精确等式，使最优解可经一次线性扫描定位。

\subsection{最优分割连续性引理}
\textbf{引理1（最优分割的连续性）}\quad 按$u_i$值升序排列后，存在最优解$\mathbf{d}^{*}$，使得分割在排序中是连续的，即存在$p \in \{0, 1, \ldots, k\}$，满足

\begin{equation}
    d_{(i)}^{*} = L, \;\forall\, i \leq p; \quad d_{(i)}^{*} = R, \;\forall\, i > p
    \label{eq:continuous_partition}
\end{equation}

\textbf{证明}\quad 反证法。设某个最优解$\mathbf{d}^{*}$中存在排序位置$a < b < c$，使得$d_{(a)} = R$，$d_{(b)} = L$，$d_{(c)} = R$。构造新解$\mathbf{d}'$，将$d_{(b)}$从$L$改为$R$，其余不变。

分析上界变化，令$\mathcal{R}' = \mathcal{R} \cup \{(b)\}$。由于$(a) \in \mathcal{R}$且按排序有$u_{(a)} \leq u_{(b)}$，故

\begin{equation}
    l^{+}(\mathbf{d}^{*}) \leq u_{(a)} \leq u_{(b)}
\end{equation}

因此$\min_{i \in \mathcal{R}'} u_i = \min_{i \in \mathcal{R}} u_i$，即$l^{+}(\mathbf{d}') = l^{+}(\mathbf{d}^{*})$，上界不变。

分析下界变化，令$\mathcal{L}' = \mathcal{L} \setminus \{(b)\}$。集合缩小，取$\max$的结果不增，故下界不升。若$\mathcal{L}' = \emptyset$，按约定$\max$取道路左边界形成的平凡下界$l_{road}^{-}$，下界不升的结论仍成立。

\begin{equation}
    l^{-}(\mathbf{d}') = \max_{i \in \mathcal{L}'} v_i \leq \max_{i \in \mathcal{L}} v_i = l^{-}(\mathbf{d}^{*})
\end{equation}



综合两方面，$W(\mathbf{d}') = l^{+}(\mathbf{d}') - l^{-}(\mathbf{d}') \geq l^{+}(\mathbf{d}^{*}) - l^{-}(\mathbf{d}^{*}) = W(\mathbf{d}^{*})$。因此$\mathbf{d}'$也是最优解，且消除了$(a), (b), (c)$处的一个不连续点。对所有不连续点重复此操作，有限次后即可得到满足式\eqref{eq:continuous_partition}的连续分割。

\textbf{推论}\quad 搜索空间从$2^k$降为$k+1$个候选分割点$p = 0, 1, \ldots, k$。其中$p = 0$表示所有障碍物均从右侧绕行即自车走最右侧通道，$p = k$表示所有障碍物均从左侧绕行即自车走最左侧通道，$1 \leq p \leq k-1$表示自车从第$p$和第$p+1$个障碍物之间的间隙通过。

\subsection{最大间隙最优性定理}
基于引理1，只需枚举$k+1$个候选分割点。对于分割点$p$，定义左绕集合$\mathcal{L}_p = \{(1), \ldots, (p)\}$和右绕集合$\mathcal{R}_p = \{(p+1), \ldots, (k)\}$。

定义$k+1$个间隙，包含道路边界形成的两个端部间隙

\begin{equation}
    g_p = \begin{cases}
        u_{(1)} - l_{road}^{-} & p = 0 \\[0.3em]
        u_{(p+1)} - v_{(p)} & p = 1, 2, \ldots, k-1 \\[0.3em]
        l_{road}^{+} - v_{(k)} & p = k
    \end{cases}
    \label{eq:gap_def}
\end{equation}

$g_0$表示道路右边界到第一个即最右侧障碍物之间的间隙，$g_k$表示最后一个即最左侧障碍物到道路左边界之间的间隙，$g_p$，其中$1 \leq p \leq k-1$，表示排序中相邻障碍物之间的间隙。

\textbf{定理1（最大间隙最优性）}\quad 对于分割点$p$，通行带宽度$W(p) = g_p$。因此最优通行带宽度为

\begin{equation}
    W^{*} = \max_{p=0,1,\ldots,k} g_p
    \label{eq:max_gap_opt}
\end{equation}

\textbf{证明}\quad 对各分割点$p$直接代入验证。

当$p = 0$即所有障碍物均为右绕时，$\mathcal{L}_0 = \emptyset$，$\mathcal{R}_0 = \{(1), \ldots, (k)\}$

\begin{equation}
    l^{+}(0) = \min(l_{road}^{+},\; u_{(1)}) = u_{(1)}, \quad l^{-}(0) = l_{road}^{-}
\end{equation}

故$W(0) = u_{(1)} - l_{road}^{-} = g_0$。

当$1 \leq p \leq k-1$时，$\mathcal{L}_p = \{(1), \ldots, (p)\}$，$\mathcal{R}_p = \{(p+1), \ldots, (k)\}$

\begin{equation}
    l^{+}(p) = \min(l_{road}^{+},\; u_{(p+1)}) = u_{(p+1)}, \quad l^{-}(p) = \max(l_{road}^{-},\; v_{(p)}) = v_{(p)}
\end{equation}

故$W(p) = u_{(p+1)} - v_{(p)} = g_p$。

当$p = k$即所有障碍物均为左绕时，$\mathcal{L}_k = \{(1), \ldots, (k)\}$，$\mathcal{R}_k = \emptyset$

\begin{equation}
    l^{+}(k) = l_{road}^{+}, \quad l^{-}(k) = \max(l_{road}^{-},\; v_{(k)}) = v_{(k)}
\end{equation}

故$W(k) = l_{road}^{+} - v_{(k)} = g_k$。

以上$l^{+}$和$l^{-}$的化简中，对中间分割点$1 \leq p \leq k-1$假设$l_{road}^{+} \geq u_{(p+1)}$且$l_{road}^{-} \leq v_{(p)}$。该假设在障碍物位于道路范围内时自然成立，因为$u_{(p+1)} = l_{(p+1)}^{-} - \delta_{(p+1)} < l_{road}^{+}$，$v_{(p)} = l_{(p)}^{+} + \delta_{(p)} > l_{road}^{-}$。对端部分割点$p = 0$或$p = k$，道路边界可能构成约束瓶颈。分以下两种情形讨论。其一，若对所有障碍物均有$u_{(1)} \leq l_{road}^{+}$且$v_{(k)} \geq l_{road}^{-}$即所有障碍物均在道路范围内，则$W(p) = g_p$对$p = 0, 1, \ldots, k$均严格成立。其二，若存在$u_{(1)} > l_{road}^{+}$，说明该障碍物已完全超出道路左边界，不构成右绕约束，几何上应将其从考虑集合中剔除，此时$k$减$1$、间隙重新编号后回到情形一。$v_{(k)} < l_{road}^{-}$的情形对称处理。因此在剔除道路边界外的障碍物后，选取最大$g_p$对应的分割点$p^{*}$即为最大化$W(p)$的分割点，$W^{*} = \max_p g_p$严格成立。

图\ref{fig:max_gap_theorem}以一维几何图示给出引理1与定理1的直观证明。上半部分将排序后的$k=4$个障碍物投影在$l$轴上，每个障碍物两侧以虚线标出各自的安全边界$u_i$与$v_i$。相邻障碍物之间以及障碍物与道路边界之间形成$k+1=5$个候选间隙$g_0,\ldots,g_4$。粉色加粗的$g_2^\star$即为该算例中的最大间隙，对应最优分割点$p^\star=2$。下半部分给出该分割对应的决策向量$\mathbf{d}^\star=(R,R,L,L)$，右绕群$\mathcal{R}$与左绕群$\mathcal{L}$在横向排序中形成连续区段，没有任何“RLR”式的交错，这正是引理1所保证的最优分割连续性。可见，原本$2^k=16$种组合中的最优解，必然落在$k+1=5$个连续分割对应的候选解之中，算法只需线性扫描比较间隙大小即可定位。

\begin{figure}[htbp]
\centering
\resizebox{0.8\textwidth}{!}{% 引理1+定理1 几何图证（1D 横向数轴）
\begin{tikzpicture}[scale=1.0, transform shape]

% 障碍物位置设计: g₂ 为最大间隙
% O₁@1.5, O₂@3.0, O₃@6.5, O₄@8.0 (中心)
% 半宽 0.35, δ 扩展 0.55
% l_road⁻=0, l_road⁺=9.5

% === 上图: 排序后的障碍物横向投影 ===

% 道路区域淡色背景（先画，避免盖住坐标轴）
\fill[bglight] (0,-0.7) rectangle (9.5,0.7);

% l 轴
\draw[->, thick] (-0.5,0) -- (10.2,0) node[right,font=\footnotesize] {$l$};

% 道路边界（垂直于 l 轴的竖线）
\draw[deepblue, very thick] (0,-0.9) -- (0,0.9);
\node[deepblue, font=\scriptsize, below] at (0,-0.95) {$l_{road}^-$};
\draw[deepblue, very thick] (9.5,-0.9) -- (9.5,0.9);
\node[deepblue, font=\scriptsize, below] at (9.5,-0.95) {$l_{road}^+$};

% 4 个障碍物（按 u_i 排序）
\fill[obsstatic] (1.15,-0.35) rectangle (1.85,0.35);
\node[white, font=\scriptsize\bfseries] at (1.5,0) {$O_1$};

\fill[obsstatic] (2.65,-0.35) rectangle (3.35,0.35);
\node[white, font=\scriptsize\bfseries] at (3.0,0) {$O_2$};

\fill[obsstatic] (6.15,-0.35) rectangle (6.85,0.35);
\node[white, font=\scriptsize\bfseries] at (6.5,0) {$O_3$};

\fill[obsstatic] (7.65,-0.35) rectangle (8.35,0.35);
\node[white, font=\scriptsize\bfseries] at (8.0,0) {$O_4$};

% u_i / v_i 膨胀区
\draw[obsbuf, dashed] (0.95,-0.55) rectangle (2.05,0.55);
\draw[obsbuf, dashed] (2.45,-0.55) rectangle (3.55,0.55);
\draw[obsbuf, dashed] (5.95,-0.55) rectangle (7.05,0.55);
\draw[obsbuf, dashed] (7.45,-0.55) rectangle (8.55,0.55);

% u_i / v_i 标注（选择性标注避免拥挤）
\node[deepblue, font=\tiny] at (0.95,0.75) {$u_1$};
\node[deeppink, font=\tiny] at (2.05,0.75) {$v_1$};
\node[deepblue, font=\tiny] at (2.45,0.75) {$u_2$};
\node[deeppink, font=\tiny] at (3.55,0.75) {$v_2$};
\node[deepblue, font=\tiny] at (5.95,0.75) {$u_3$};
\node[deeppink, font=\tiny] at (7.05,0.75) {$v_3$};
\node[deepblue, font=\tiny] at (7.45,0.75) {$u_4$};
\node[deeppink, font=\tiny] at (8.55,0.75) {$v_4$};

% 5 个间隙标注（在下方）
% g₀ = u₁ - l_road⁻ = 0.95
\draw[<->, thick, lightgreen!60!black] (0,-1.2) -- (0.95,-1.2);
\node[font=\scriptsize, lightgreen!60!black] at (0.475,-1.5) {$g_0$};

% g₁ = u₂ - v₁ = 0.40
\draw[<->, thick, lightgreen!60!black] (2.05,-1.2) -- (2.45,-1.2);
\node[font=\scriptsize, lightgreen!60!black] at (2.25,-1.5) {$g_1$};

% g₂ = u₃ - v₂ = 2.40 ← 最大间隙
\draw[<->, very thick, deeppink] (3.55,-1.2) -- (5.95,-1.2);
\node[font=\small\bfseries, deeppink] at (4.75,-1.5) {$g_2^\star$};

% g₃ = u₄ - v₃ = 0.40
\draw[<->, thick, lightgreen!60!black] (7.05,-1.2) -- (7.45,-1.2);
\node[font=\scriptsize, lightgreen!60!black] at (7.25,-1.5) {$g_3$};

% g₄ = l_road⁺ - v₄ = 0.95
\draw[<->, thick, lightgreen!60!black] (8.55,-1.2) -- (9.5,-1.2);
\node[font=\scriptsize, lightgreen!60!black] at (9.025,-1.5) {$g_4$};

% 最优分割线
\draw[deeppink, thick, dashed] (4.75,-0.9) -- (4.75,0.9);

\node[font=\scriptsize] at (4.75,-1.9) {$k{=}4$ 个障碍物，$k{+}1{=}5$ 个候选间隙，$g_2^\star$ 最大};

% === 下图: 决策向量 ===
\begin{scope}[shift={(0,-3.8)}]

  % R 群背景
  \fill[lightgreen!20] (0,-0.6) rectangle (4.75,0.6);
  % L 群背景
  \fill[improvedpink!25] (4.75,-0.6) rectangle (9.5,0.6);

  % 4 个障碍物
  \fill[obsstatic] (1.15,-0.35) rectangle (1.85,0.35);
  \node[white, font=\scriptsize\bfseries] at (1.5,0) {$O_1$};
  \node[deepblue, font=\scriptsize\bfseries] at (1.5,-0.65) {$d_1{=}R$};

  \fill[obsstatic] (2.65,-0.35) rectangle (3.35,0.35);
  \node[white, font=\scriptsize\bfseries] at (3.0,0) {$O_2$};
  \node[deepblue, font=\scriptsize\bfseries] at (3.0,-0.65) {$d_2{=}R$};

  \fill[obsstatic] (6.15,-0.35) rectangle (6.85,0.35);
  \node[white, font=\scriptsize\bfseries] at (6.5,0) {$O_3$};
  \node[deeppink, font=\scriptsize\bfseries] at (6.5,-0.65) {$d_3{=}L$};

  \fill[obsstatic] (7.65,-0.35) rectangle (8.35,0.35);
  \node[white, font=\scriptsize\bfseries] at (8.0,0) {$O_4$};
  \node[deeppink, font=\scriptsize\bfseries] at (8.0,-0.65) {$d_4{=}L$};

  % 分割线
  \draw[deeppink, thick, dashed] (4.75,-0.8) -- (4.75,0.8);
  \node[deeppink, font=\scriptsize] at (4.75,1.0) {$p^\star{=}2$};

  % 群标注
  \node[deepblue, font=\scriptsize] at (2.25,-1.1) {$\mathcal{R}{=}\{O_1,O_2\}$};
  \node[deeppink, font=\scriptsize] at (7.25,-1.1) {$\mathcal{L}{=}\{O_3,O_4\}$};

  % 连续性说明
  \node[font=\scriptsize] at (4.75,-1.5) {$\mathbf{d}^\star{=}(R,R,L,L)$，排序中连续，无交错};
\end{scope}

\end{tikzpicture}
}
\caption{最优分割连续性与最大间隙最优性的几何图证}
\label{fig:max_gap_theorem}
\end{figure}

至此，最优通行带等价于最宽横向间隙的选取。引理1的交换论证排除了非连续分配更优的可能性，定理1将通行带宽度与间隙值建立精确等式，使最优解可经一次线性扫描获得。

几个边界情形需单独处理。第一，当某些障碍物在横向上存在重叠即$v_{(p)} > u_{(p+1)}$时，对应的间隙$g_p$为负值，表示该通道不可通行，算法自然跳过。第二，当所有$k+1$个间隙均为负或不足车宽时，$W^{*} \leq 0$，说明本车道确实不存在可通行路径，这是真实的BLOCKED而非算法的局限。第三，引理1的证明仅要求$u$值按排序单调不减，当存在$u_i$相同的障碍物时，可按$v_i$的降序作为二级排序键，不影响引理和定理的成立。

\section{算法描述与复杂度分析}
算法\ref{alg:optimal_band}给出了多障碍物统一规划求解的完整流程，将引理1和定理1转化为可执行的计算步骤，依次完成到达时间估计、时变SL投影、扫描线分组、间隙计算与最大间隙选取、避让方向分配、路径边界构建和速度约束输出七个步骤。总体结构如图\ref{fig:algorithm_flow}所示：从参考线与障碍物集合起步，先经威胁度评估与动态安全距离确定每个障碍物的$\delta_i$，再由扫描线完成纵向重叠分组，组内按$u_i$排序查找最大间隙并分配避让方向，最后将所有组的边界合并为全局\texttt{path\_boundary}并投影至ST空间生成走廊。图中每一步骤右侧标注了复杂度，主要开销集中在扫描线分组的$O(n\log n)$，其余步骤均为线性，总体满足$O(n\log n + n N_s)$的实时性目标。

\begin{figure}[htbp]
\centering
\resizebox{0.8\textwidth}{!}{% MIKU 多障碍物统一规划算法总流程
\begin{tikzpicture}[
    scale=1.0, transform shape,
    stepimp/.style={rectangle, rounded corners=3pt, draw=deeppink, fill=improvedpink, minimum width=5.4cm, minimum height=1.0cm, font=\small, align=center},
    inp/.style={rectangle, draw=deepblue, fill=inputyellow, minimum width=4.0cm, minimum height=0.8cm, font=\footnotesize, align=center},
    outnode/.style={rectangle, rounded corners=3pt, draw=deeppink, fill=outputpink, minimum width=4.0cm, minimum height=0.8cm, font=\small\bfseries, align=center},
    qpnode/.style={rectangle, rounded corners=3pt, draw=deepblue, fill=lightblue, minimum width=4.0cm, minimum height=0.8cm, font=\small\bfseries, align=center},
    arr/.style={-Stealth, thick, deepblue},
    injarr/.style={-Stealth, thick, deeppink},
]

% 输入
\node[inp] (i1) at (-5.8, 6.0) {参考线 $+$ 障碍物集合};
\node[inp] (i2) at (-5.8, 5.0) {预测轨迹 $+$ 车辆状态};

% 步骤（间距 1.7，全部为 MIKU 改进）
\node[stepimp] (s1) at (0, 5.5)  {\textbf{Step 1} 到达时间 $\tau(s)$ 与时变 SL 投影};
\node[stepimp] (s2) at (0, 3.8)  {\textbf{Step 2} 多因子威胁度评估 $\Theta_i$};
\node[stepimp] (s3) at (0, 2.1)  {\textbf{Step 3} 动态安全裕度 $\delta_i$ 与边界 $u_i,v_i$};
\node[stepimp] (s4) at (0, 0.4)  {\textbf{Step 4} 扫描线重叠分组 $\mathcal{G}_1,\ldots,\mathcal{G}_m$};
\node[stepimp] (s5) at (0,-1.3)  {\textbf{Step 5} 组内最大间隙 $g^\star_j=\max_k g_{j,k}$};
\node[stepimp] (s6) at (0,-3.0)  {\textbf{Step 6} 合并全局路径横向边界};
\node[stepimp] (s7) at (0,-4.7)  {\textbf{Step 7} ST 时空可通行走廊 $\Omega$};

% 输出（左侧）与跨接口的速度 QP
\node[outnode] (o1) at (-5.8,-3.0) {路径横向边界};
\node[outnode] (o2) at (-5.8,-4.7) {时空可通行走廊};
\node[qpnode]  (qp) at (-5.8,-6.4) {速度 QP（求解器不变）};

% 输入箭头
\draw[arr] (i1.east) -- (s1.west);
\draw[arr] (i2.east) -- (s1.west);

% 步骤间箭头
\draw[arr] (s1) -- (s2);
\draw[arr] (s2) -- (s3);
\draw[arr] (s3) -- (s4);
\draw[arr] (s4) -- (s5);
\draw[arr] (s5) -- (s6);
\draw[arr] (s6) -- (s7);

% 输出箭头
\draw[arr] (s6.west) -- (o1.east);
\draw[arr] (s7.west) -- (o2.east);

% 走廊注入速度 QP（跨解耦接口）
\draw[injarr] (o2.south) -- node[right, font=\footnotesize, deeppink] {逐时间步纵向上界} (qp.north);

\end{tikzpicture}
}
\caption{多障碍物统一规划算法的总体流程与复杂度}
\label{fig:algorithm_flow}
\end{figure}

\begin{algorithm}[htbp]
\caption{多障碍物统一规划求解算法}
\label{alg:optimal_band}
\KwIn{$n$个障碍物，含预测轨迹，道路边界$[l_{road}^{-}(s), l_{road}^{+}(s)]$、自车状态$(s_{ego}, v_{ego}, a_{ego})$、各障碍物安全裕度$\delta_i$}
\KwOut{路径边界$\{[l^{-}(s), l^{+}(s)]\}$、速度约束集合$\mathcal{T}$}
\textcolor{gray!30!black}{\textit{//步骤1：到达时间估计}}\;
\ForEach{离散位置 $s$}{
    \eIf{$|a_{ego}| > \epsilon$}{
        $\tau(s) \leftarrow \left(-v_{ego} + \sqrt{v_{ego}^2 + 2\,a_{ego}\,(s - s_{ego})}\right) / a_{ego}$ \;
    }{
        $\tau(s) \leftarrow (s - s_{ego}) / v_{ego}$ \;
    }
}
\textcolor{gray!30!black}{\textit{//步骤2：时变SL投影}}\;
\ForEach{障碍物 $O_i$}{
    查询预测轨迹，计算$O_i(\tau(s))$ \;
    \If{$O_i$在$\tau(s)$时刻已移出参考线附近}{
        标记$O_i$在位置$s$处不活跃 \;
    }
    $u_i(s) \leftarrow l_i^{-}(\tau(s)) - \delta_i$，$v_i(s) \leftarrow l_i^{+}(\tau(s)) + \delta_i$ \;
}
\textcolor{gray!30!black}{\textit{//步骤3：按$s_i^{-}$排序，扫描线分组}}\;
将所有障碍物按$s_i^{-}$升序排序 \;
$s_{max} \leftarrow -\infty$，$group \leftarrow \emptyset$ \;
\ForEach{排序后的障碍物 $O_i$}{
    \eIf{$s_i^{-} \leq s_{max}$}{
        将$O_i$加入当前分量，$s_{max} \leftarrow \max(s_{max}, s_i^{+})$ \;
    }{
        保存当前分量，开始新分量$\{O_i\}$，$s_{max} \leftarrow s_i^{+}$ \;
    }
}
\textcolor{gray!30!black}{\textit{//步骤4：组内按$u_i$排序，计算间隙}}\;
\ForEach{连通分量 $C_j$，含$k$个障碍物}{
    按$u_i$升序排列为$O_{(1)}, \ldots, O_{(k)}$ \;
    计算$k+1$个间隙$g_0, g_1, \ldots, g_k$（式\eqref{eq:gap_def}） \;
    \textcolor{gray!30!black}{\textit{//步骤5：选最大间隙，分配避让方向}}\;
    $p^{*} \leftarrow \arg\max_{p} g_p$ \;
    $d_{(i)} \leftarrow L$，$\forall\, i \leq p^{*}$；$d_{(i)} \leftarrow R$，$\forall\, i > p^{*}$ \;
}
\textcolor{gray!30!black}{\textit{//步骤6：构建路径边界}}\;
\ForEach{离散位置 $s$}{
    按分配的避让方向更新$[l^{-}(s), l^{+}(s)]$ \;
}
\textcolor{gray!30!black}{\textit{//步骤7：输出速度约束}}\;
\ForEach{依赖动态障碍物移开的位置 $s_k$}{
    $\mathcal{T} \leftarrow \mathcal{T} \cup \{(s_k, \tau(s_k))\}$ \;
}
\KwRet $\{[l^{-}(s), l^{+}(s)]\}$，$\mathcal{T}$ \;
\end{algorithm}

下面分析算法\ref{alg:optimal_band}各步骤的时间和空间复杂度。

对单个含$k$个障碍物的连通分量，算法步骤4中按$u_i$排序的时间复杂度为$O(k \log k)$，计算$k+1$个间隙和选取最大值均为$O(k)$，因此单分量的时间复杂度为$O(k \log k)$。空间复杂度为$O(k)$，用于存储排序后的障碍物列表和间隙数组。

全局来看，步骤1和步骤2对每个障碍物执行常数时间操作，复杂度为$O(n)$。步骤3中按$s_i^{-}$排序的复杂度为$O(n \log n)$，扫描线遍历为$O(n)$。步骤4--5中，所有连通分量的组内排序总复杂度满足$\sum_j O(k_j \log k_j) \leq O(n \log n)$，其中$\sum_j k_j = n$，由Jensen不等式的推广可得。步骤6和步骤7均为$O(K)$，其中$K$为离散位置数。因此，算法的总时间复杂度为$O(n \log n + K)$，其中$O(n \log n)$为障碍物处理，$O(K)$为路径边界构建。在$n \ll K$的典型场景下，$O(K)$项占主导。

表\ref{tab:algo_compare}从多个维度对比了本文算法与Apollo Baseline的差异。

\begin{table}[htbp]
\centering
\caption{本文算法与Apollo Baseline的对比}
\label{tab:algo_compare}
\small
\begin{tabularx}{\textwidth}{lXX}
\toprule
& \textbf{Apollo Baseline} & \textbf{本文方法} \\
\midrule
\textbf{决策粒度} & 逐障碍物独立决策 & 组级全局决策 \\
\textbf{最优性保证} & 贪心策略，无最优性保证 & 定理保证组内最优 \\
\textbf{时间信息} & path阶段仅$t=0$快照（speed阶段使用预测轨迹） & 预测轨迹$O_i(\tau(s))$ \\
\textbf{动态障碍物} & 不参与路径决策 & 通过预测位置参与 \\
\textbf{BLOCKED处理} & 截断路径$\to$停车 & 时变间隙可能打开$\to$继续通行 \\
\textbf{速度协调} & 无 & 输出$\tau(s_k)$约束 \\
\textbf{时间复杂度} & $O(nK)$ & $O(n \log n)$ \\
\bottomrule
\end{tabularx}
\end{table}

Baseline在每个离散位置$s$处遍历所有障碍物共$K$个离散点，总复杂度为$O(nK)$；本文算法在排序后只需对$k+1$个候选分割点做线性扫描，复杂度为$O(n\log n)$。两者在常规多障碍物规模下均可在毫秒级完成，远低于100\,ms规划周期限制。本文算法的核心价值在于保证决策的全局最优性，而非计算效率本身的提升。

\section{数值仿真}
\label{sec:exp_ch6}

本节通过两个对照场景验证扫描线分组与最大间隙策略对Baseline独立决策矛盾的修正效果。场景二为动+静混合的小规模分组，场景三为单车道维修封闭+借道的大规模连通分量，分别检验组内方向冲突的破解与跨车道全局协调的能力。

\subsection{场景二动静混合的组内方向协调}

场景二由一名横穿行人加左前方静态停车构成，行人与停车在纵向上存在重叠区，构成同一连通分量。Baseline逐障碍物决策时，行人横向位置触发右推，停车横向位置又强制左推，两者冲突使$l_{\min}>l_{\max}$。MIKU将两者并入同一组，按$u_i$排序计算两个候选间隙，选取最大间隙$g^{*}$对应的统一避让方向。

\begin{figure}[H]
\centering
% 场景02 SL 平面 Baseline vs MIKU 对比（上下排，每子图占满 \textwidth）
\def\datapath{data/02_ped_plus_parked/}
\pgfplotsset{
    discard if not/.style 2 args={x filter/.code={
        \edef\tempa{\thisrow{#1}}\edef\tempb{#2}
        \ifx\tempa\tempb\else\def\pgfmathresult{nan}\fi
    }},
    discard if not2/.style n args={4}{x filter/.code={
        \edef\tempa{\thisrow{#1}}\edef\tempb{#2}
        \edef\tempc{\thisrow{#3}}\edef\tempd{#4}
        \ifx\tempa\tempb
            \ifx\tempc\tempd\else\def\pgfmathresult{nan}\fi
        \else\def\pgfmathresult{nan}\fi
    }}
}
\begin{subfigure}[b]{\textwidth}
\centering
\begin{tikzpicture}
\begin{axis}[
    width=\linewidth, height=4.2cm,
    xlabel={$s$ (m)}, ylabel={$l$ (m)},
    xmin=0, xmax=32.0, ymin=-2.2, ymax=2.2,
    grid=both, grid style={gray!25},
    axis line style={thick},
    tick label style={font=\scriptsize},
    label style={font=\footnotesize},
    legend style={at={(0.98,0.02)}, anchor=south east, font=\scriptsize, draw=none, fill=white, fill opacity=0.85},
]
\addplot[fill=bglight, draw=none, forget plot, on layer=axis background] coordinates {(0,-1.875) (32.0,-1.875) (32.0,1.875) (0,1.875)} \closedcycle;

\draw[black!55, semithick] (axis cs:0,-1.875) -- (axis cs:32.0,-1.875);
\draw[black!55, semithick] (axis cs:0,1.875) -- (axis cs:32.0,1.875);
\draw[black!30, dashed, thin] (axis cs:0,0) -- (axis cs:32.0,0);
% 可行带阴影（blocked 时按 blocked_s 截断）
\addplot[name path=lmin_baseline, draw=vibrantorange, thin, dashed, forget plot, ]
    table[col sep=comma, x=s, y=l_min, discard if not={mode}{baseline}] {\datapath sl.csv};
\addplot[name path=lmax_baseline, draw=vibrantorange, thin, dashed, forget plot, ]
    table[col sep=comma, x=s, y=l_max, discard if not={mode}{baseline}] {\datapath sl.csv};
\addplot[fill=lightgreen!35, draw=none, forget plot] fill between[of=lmin_baseline and lmax_baseline];
% 路径（blocked 时按 blocked_s 截断）
\addplot[deepblue, line width=2pt, unbounded coords=discard] table[col sep=comma, x=s, y=l_path, discard if not={mode}{baseline}] {\datapath sl.csv};
\addlegendentry{Baseline}


% 障碍物 行人
\fill[dangerred, opacity=0.95] (axis cs:10.000,-0.500) circle [radius=5pt];
\fill[dangerred, opacity=0.55] (axis cs:10.000,0.292) circle [radius=5pt];
\fill[dangerred, opacity=0.35] (axis cs:10.000,1.084) circle [radius=5pt];
\fill[dangerred, opacity=0.20] (axis cs:10.000,1.900) circle [radius=5pt];
\draw[->, dangerred, very thick, opacity=0.85] (axis cs:10.000,-0.500) -- (axis cs:10.000,1.900);
\node[anchor=south, font=\tiny, color=dangerred] at (axis cs:10.000,0.250) {行人 $v=(+0.0,+1.2)$};
% 障碍物 停车
\fill[dangerred!70, draw=dangerred, thick] (axis cs:20.000,0.800) rectangle (axis cs:24.000,1.800);
\node[anchor=south, font=\tiny, color=dangerred!80!black] at (axis cs:22.000,2.100) {停车};
\node[rectangle, fill=deepblue!75, draw=deepblue, thick, minimum width=18pt, minimum height=8pt, inner sep=0pt] (egobox) at (axis cs:0.000,0.000) {};
\fill[white] ([xshift=0pt]egobox.east) -- ([xshift=-3pt,yshift=2.5pt]egobox.east) -- ([xshift=-3pt,yshift=-2.5pt]egobox.east) -- cycle;
\node[anchor=south west, font=\tiny\bfseries, color=deepblue] at (axis cs:0.20,0.95) {EGO};
\end{axis}
\end{tikzpicture}
\caption{Baseline}
\end{subfigure}

\vspace{0.6em}

\begin{subfigure}[b]{\textwidth}
\centering
\begin{tikzpicture}
\begin{axis}[
    width=\linewidth, height=4.2cm,
    xlabel={$s$ (m)}, ylabel={$l$ (m)},
    xmin=0, xmax=32.0, ymin=-2.2, ymax=2.2,
    grid=both, grid style={gray!25},
    axis line style={thick},
    tick label style={font=\scriptsize},
    label style={font=\footnotesize},
    legend style={at={(0.98,0.02)}, anchor=south east, font=\scriptsize, draw=none, fill=white, fill opacity=0.85},
]
\addplot[fill=bglight, draw=none, forget plot, on layer=axis background] coordinates {(0,-1.875) (32.0,-1.875) (32.0,1.875) (0,1.875)} \closedcycle;

\draw[black!55, semithick] (axis cs:0,-1.875) -- (axis cs:32.0,-1.875);
\draw[black!55, semithick] (axis cs:0,1.875) -- (axis cs:32.0,1.875);
\draw[black!30, dashed, thin] (axis cs:0,0) -- (axis cs:32.0,0);
% 可行带阴影（blocked 时按 blocked_s 截断）
\addplot[name path=lmin_miku, draw=vibrantorange, thin, dashed, forget plot, ]
    table[col sep=comma, x=s, y=l_min, discard if not={mode}{miku}] {\datapath sl.csv};
\addplot[name path=lmax_miku, draw=vibrantorange, thin, dashed, forget plot, ]
    table[col sep=comma, x=s, y=l_max, discard if not={mode}{miku}] {\datapath sl.csv};
\addplot[fill=lightgreen!35, draw=none, forget plot] fill between[of=lmin_miku and lmax_miku];
% 路径（blocked 时按 blocked_s 截断）
\addplot[vibrantgreen, line width=2pt, unbounded coords=discard] table[col sep=comma, x=s, y=l_path, discard if not={mode}{miku}] {\datapath sl.csv};
\addlegendentry{MIKU}


% 障碍物 行人
\fill[dangerred, opacity=0.95] (axis cs:10.000,-0.500) circle [radius=5pt];
\fill[dangerred, opacity=0.55] (axis cs:10.000,0.292) circle [radius=5pt];
\fill[dangerred, opacity=0.35] (axis cs:10.000,1.084) circle [radius=5pt];
\fill[dangerred, opacity=0.20] (axis cs:10.000,1.900) circle [radius=5pt];
\draw[->, dangerred, very thick, opacity=0.85] (axis cs:10.000,-0.500) -- (axis cs:10.000,1.900);
\node[anchor=south, font=\tiny, color=dangerred] at (axis cs:10.000,0.250) {行人 $v=(+0.0,+1.2)$};
% 障碍物 停车
\fill[dangerred!70, draw=dangerred, thick] (axis cs:20.000,0.800) rectangle (axis cs:24.000,1.800);
\node[anchor=south, font=\tiny, color=dangerred!80!black] at (axis cs:22.000,2.100) {停车};
\node[rectangle, fill=deepblue!75, draw=deepblue, thick, minimum width=18pt, minimum height=8pt, inner sep=0pt] (egobox) at (axis cs:0.000,0.000) {};
\fill[white] ([xshift=0pt]egobox.east) -- ([xshift=-3pt,yshift=2.5pt]egobox.east) -- ([xshift=-3pt,yshift=-2.5pt]egobox.east) -- cycle;
\node[anchor=south west, font=\tiny\bfseries, color=deepblue] at (axis cs:0.20,0.95) {EGO};
\end{axis}
\end{tikzpicture}
\caption{MIKU}
\end{subfigure}

\caption{场景二SL平面路径与边界对比}
\label{fig:exp_ch6_sl_2}
\end{figure}

\begin{figure}[H]
\centering
\input{../图片/fig_exp_02_ped_plus_parked_st}
\caption{场景二ST平面速度规划对比}
\label{fig:exp_ch6_st_2}
\end{figure}

\begin{figure}[H]
\centering
% 场景02 速度/加速度 Baseline vs GTOC 对比
\def\datapath{data/02_ped_plus_parked/}
\pgfplotsset{
    discard if not/.style 2 args={x filter/.code={
        \edef\tempa{\thisrow{#1}}\edef\tempb{#2}
        \ifx\tempa\tempb\else\def\pgfmathresult{nan}\fi
    }},
    discard if not2/.style n args={4}{x filter/.code={
        \edef\tempa{\thisrow{#1}}\edef\tempb{#2}
        \edef\tempc{\thisrow{#3}}\edef\tempd{#4}
        \ifx\tempa\tempb
            \ifx\tempc\tempd\else\def\pgfmathresult{nan}\fi
        \else\def\pgfmathresult{nan}\fi
    }}
}
\begin{subfigure}[b]{\textwidth}
\centering
\begin{tikzpicture}
\begin{axis}[
    width=\linewidth, height=4.2cm,
    xlabel={$t$ (s)}, ylabel={$v$ (m/s)},
    xmin=0, xmax=6.5, ymin=0, ymax=10.4,
    grid=both, grid style={gray!25},
    axis line style={thick},
    tick label style={font=\tiny},
    label style={font=\scriptsize},
    legend style={at={(0.5,-0.34)}, anchor=north, legend columns=2,
        font=\fontsize{6}{7}\selectfont, draw=none, /tikz/every even column/.append style={column sep=2em}},
]
% v_ref 参考线
\draw[black!50, dashed, thick] (axis cs:0,8.00) -- (axis cs:6.50,8.00);
\node[black!50, font=\tiny, anchor=south east] at (axis cs:6.50,8.00) {$v_{ref}$};
\addplot[deepblue, line width=2pt]
    table[col sep=comma, x=t, y=v_qp, discard if not={mode}{baseline}] {\datapath st_curves.csv};
\addlegendentry{Baseline\,(\,$\bar{v}{=}4.11$,\,$|a|_{max}{=}4.00$,\,$|j|_{max}{=}13.31$,\,$s_{end}{=}26.4$\,m\,)}
\addplot[vibrantgreen, line width=2pt]
    table[col sep=comma, x=t, y=v_qp, discard if not={mode}{gtoc}] {\datapath st_curves.csv};
\addlegendentry{GTOC\,(\,$\bar{v}{=}8.00$,\,$|a|_{max}{=}0.00$,\,$|j|_{max}{=}0.00$,\,$s_{end}{=}52.0$\,m\,)}
\end{axis}
\end{tikzpicture}
\caption{速度对比}
\end{subfigure}

\vspace{0.8em}

\begin{subfigure}[b]{\textwidth}
\centering
\begin{tikzpicture}
\begin{axis}[
    width=\linewidth, height=4.2cm,
    xlabel={$t$ (s)}, ylabel={$a$ (m/s$^2$)},
    xmin=0, xmax=6.5, ymin=-5, ymax=5,
    restrict y to domain*=-5:5,
    unbounded coords=jump,
    grid=both, grid style={gray!25},
    axis line style={thick},
    tick label style={font=\tiny},
    label style={font=\scriptsize},
    legend style={at={(0.5,-0.34)}, anchor=north, legend columns=4,
        font=\fontsize{6}{7}\selectfont, draw=none, /tikz/every even column/.append style={column sep=1em}},
]
% Baseline a_x 实线
\addplot[deepblue, line width=1.8pt]
    table[col sep=comma, x=t, y=a_qp, discard if not={mode}{baseline}] {\datapath st_curves.csv};
\addlegendentry{Baseline\,$a_x$}
% GTOC a_x 实线
\addplot[vibrantgreen, line width=1.8pt]
    table[col sep=comma, x=t, y=a_qp, discard if not={mode}{gtoc}] {\datapath st_curves.csv};
\addlegendentry{GTOC\,$a_x$}
% Baseline a_y dashed
\addplot[deepblue, line width=1.4pt, dashed]
    table[col sep=comma, x=t, y=a_y, discard if not={mode}{baseline}] {\datapath st_curves.csv};
\addlegendentry{Baseline\,$a_y$}
% GTOC a_y dashed
\addplot[vibrantgreen, line width=1.4pt, dashed]
    table[col sep=comma, x=t, y=a_y, discard if not={mode}{gtoc}] {\datapath st_curves.csv};
\addlegendentry{GTOC\,$a_y$}
\end{axis}
\end{tikzpicture}
\caption{纵向/横向加速度对比（$a_x$ 实线、$a_y$ 虚线；蓝 Baseline、绿 GTOC）}
\end{subfigure}

\vspace{0.8em}

\begin{subfigure}[b]{\textwidth}
\centering
\begin{tikzpicture}
\begin{axis}[
    width=\linewidth, height=4.2cm,
    xlabel={$t$ (s)}, ylabel={$j$ (m/s$^3$)},
    xmin=0, xmax=6.5,
    grid=both, grid style={gray!25},
    axis line style={thick},
    tick label style={font=\tiny},
    label style={font=\scriptsize},
    legend style={at={(0.5,-0.34)}, anchor=north, legend columns=2,
        font=\fontsize{6}{7}\selectfont, draw=none, /tikz/every even column/.append style={column sep=2em}},
]
\addplot[deepblue, line width=1.6pt]
    table[col sep=comma, x=t, y=j_qp, discard if not={mode}{baseline}] {\datapath st_curves.csv};
\addlegendentry{Baseline\,$j$}
\addplot[vibrantgreen, line width=1.6pt]
    table[col sep=comma, x=t, y=j_qp, discard if not={mode}{gtoc}] {\datapath st_curves.csv};
\addlegendentry{GTOC\,$j$}
\end{axis}
\end{tikzpicture}
\caption{jerk 对比}
\end{subfigure}

\caption{场景二$v(t)$、$a(t)$、$j(t)$对比}
\label{fig:exp_ch6_va_2}
\end{figure}

\subsection{场景三单车道封闭导流的跨车道全局协调}

场景三由18个静态障碍构成，包括入口漏斗5个交通锥、维持段8段水马、出口漏斗5个交通锥，所有障碍属于同一大型连通分量，要求自车借用左侧相邻车道通过。Baseline在入口漏斗段对每锥独立选择左推或右推，与维持段水马跨越车道分界线的几何产生方向矛盾，在维持段入口处$l_{\min}>l_{\max}$，path被截断停车于$s\approx$\,\MFourBaselineSEnd\,m。MIKU将整段导流识别为单一组，组级最大间隙搜索将$p^{*}$分配至所有水马的左侧，路径在入口处即决定借入相邻车道深处，全程平顺通过。

\begin{figure}[H]
\centering
\input{../图片/fig_exp_04_dense_construction_sl}
\caption{场景三SL平面路径与边界对比}
\label{fig:exp_ch6_sl_4}
\end{figure}

\begin{figure}[H]
\centering
% 场景04 ST 平面 Baseline vs MIKU 对比
\def\datapath{data/04_dense_construction/}
\pgfplotsset{
    discard if not/.style 2 args={x filter/.code={
        \edef\tempa{\thisrow{#1}}\edef\tempb{#2}
        \ifx\tempa\tempb\else\def\pgfmathresult{nan}\fi
    }},
    discard if not2/.style n args={4}{x filter/.code={
        \edef\tempa{\thisrow{#1}}\edef\tempb{#2}
        \edef\tempc{\thisrow{#3}}\edef\tempd{#4}
        \ifx\tempa\tempb
            \ifx\tempc\tempd\else\def\pgfmathresult{nan}\fi
        \else\def\pgfmathresult{nan}\fi
    }}
}
\begin{subfigure}[b]{0.48\textwidth}
\centering
\begin{tikzpicture}
\begin{axis}[
    width=\linewidth, height=5.5cm,
    xlabel={$t$ (s)}, ylabel={$s$ (m)},
    xmin=0, xmax=15.0, ymin=0, ymax=127.5,
    restrict y to domain*=0:127.5,
    unbounded coords=jump,
    grid=both, grid style={gray!25},
    axis line style={thick},
    tick label style={font=\scriptsize},
    label style={font=\footnotesize},
    legend style={at={(0.02,0.98)}, anchor=north west, font=\scriptsize, draw=none, fill=white, fill opacity=0.9},
]
% 该模式无 ST 禁行区（路径阶段已绕开冲突，无投影）
% DP 粗解
\addplot[improvedpink, thick, dashed]
    table[col sep=comma, x=t, y=s_dp, discard if not={mode}{baseline}] {\datapath st_curves.csv};
\addlegendentry{$s_{dp}$}
% QP 精解
\addplot[deepblue, line width=2pt]
    table[col sep=comma, x=t, y=s_qp, discard if not={mode}{baseline}] {\datapath st_curves.csv};
\addlegendentry{$s_{qp}$ (Baseline)}

\end{axis}
\end{tikzpicture}
\caption{Baseline}
\end{subfigure}\hfill
\begin{subfigure}[b]{0.48\textwidth}
\centering
\begin{tikzpicture}
\begin{axis}[
    width=\linewidth, height=5.5cm,
    xlabel={$t$ (s)}, ylabel={$s$ (m)},
    xmin=0, xmax=15.0, ymin=0, ymax=127.5,
    restrict y to domain*=0:127.5,
    unbounded coords=jump,
    grid=both, grid style={gray!25},
    axis line style={thick},
    tick label style={font=\scriptsize},
    label style={font=\footnotesize},
    legend style={at={(0.02,0.98)}, anchor=north west, font=\scriptsize, draw=none, fill=white, fill opacity=0.9},
]
% 该模式无 ST 禁行区（路径阶段已绕开冲突，无投影）
% DP 粗解
\addplot[improvedpink, thick, dashed]
    table[col sep=comma, x=t, y=s_dp, discard if not={mode}{miku}] {\datapath st_curves.csv};
\addlegendentry{$s_{dp}$}
% QP 精解
\addplot[vibrantgreen, line width=2pt]
    table[col sep=comma, x=t, y=s_qp, discard if not={mode}{miku}] {\datapath st_curves.csv};
\addlegendentry{$s_{qp}$ (MIKU)}

\end{axis}
\end{tikzpicture}
\caption{MIKU}
\end{subfigure}

\caption{场景三ST平面速度规划对比}
\label{fig:exp_ch6_st_4}
\end{figure}

\begin{figure}[H]
\centering
\input{../图片/fig_exp_04_dense_construction_va}
\caption{场景三$v(t)$、$a(t)$、$j(t)$对比}
\label{fig:exp_ch6_va_4}
\end{figure}

\subsection{指标汇总}

两个场景的关键指标列于表\ref{tab:exp_ch6_metrics}。Baseline在两个场景下均因方向冲突触发BLOCKED，MIKU通过组级最优分配将通行成功率从0翻转为1。

\begin{table}[htbp]
\centering
\caption{场景二三Baseline与MIKU指标对照}
\label{tab:exp_ch6_metrics}
\small
\begin{tabular}{clccccc}
\toprule
\textbf{场景} & \textbf{模式} & $\bar{v}$\,(m/s) & $|a_x|_{\max}$\,(m/s$^2$) & $s_{end}$\,(m) & $t_{arrive}$\,(s) & 通过 \\
\midrule
\multirow{2}{*}{二} & Baseline & \MTwoBaselineAvgV  & \MTwoBaselineMaxAbsA  & \MTwoBaselineSEnd  & \MTwoBaselineTArrive  & $\times$ \\
                   & MIKU     & \MTwoMikuAvgV      & \MTwoMikuMaxAbsA      & \MTwoMikuSEnd      & \MTwoMikuTArrive      & $\checkmark$ \\
\midrule
\multirow{2}{*}{三} & Baseline & \MFourBaselineAvgV & \MFourBaselineMaxAbsA & \MFourBaselineSEnd & \MFourBaselineTArrive & $\times$ \\
                   & MIKU     & \MFourMikuAvgV     & \MFourMikuMaxAbsA     & \MFourMikuSEnd     & \MFourMikuTArrive     & $\checkmark$ \\
\bottomrule
\end{tabular}
\end{table}

实验结果与本章理论分析一致。场景二中MIKU的组级方向协调避免了Baseline逐障碍物贪心的方向冲突；场景三中MIKU的最大间隙搜索发现了Baseline贪心策略遗漏的统一借入相邻车道的全局最优解。两个场景表明，定理1所述组级最优对逐障碍物贪心的替代能切实提升通行成功率。

\section{本章小结}

本章针对Apollo PathBoundsDecider独立决策矛盾导致路径BLOCKED的问题，提出MIKU的第二项改进，即组级最优通行带与最大间隙策略。源码分析定位了\texttt{UpdatePathBoundary\-BySLPolygon}函数的逐障碍物贪心决策模式，对称双障碍物场景的几何分析则说明贪心策略可能在$2^k$种避让方向组合中错误选中不可行解。算法层面，基于扫描线的纵向重叠检测将相互关联的障碍物聚为连通分量，组内引理证明最优分割在横向排序中的连续性，定理证明每个候选分割的通行带宽度恰好等于对应的横向间隙，将组合搜索复杂度从$O(2^k)$降至$O(k\log k)$。两个证明对$\delta$是否为常数不做假设，第四章的差异化$\delta_i$可直接代入本章求解算法。本章输出的组级避让方向决策与通行带宽度，是第六章构造时空可通行走廊的SL层基础。
